\section{Requerimientos no funcionales}
Los siguientes requerimientos establecen restricciones de calidad y operación
para el prototipo, y complementan las operaciones de sucursales, premios y
clientes.

\begin{enumerate}[label=\textbf{RNF-\arabic*.}, leftmargin=*]
    \item \textbf{Interacción.} El prototipo deberá ejecutarse en una terminal y
    ofrecer un menú principal con acceso a las operaciones de las tres entidades.
    La navegación deberá permitir regresar o salir sin terminar inesperadamente.
    \item \textbf{Persistencia.} La información deberá almacenarse en archivos
    CSV separados, con encabezado, delimitador coma y codificación UTF-8. Los
    datos deberán poder recuperarse en ejecuciones posteriores.
    \item \textbf{Integridad.} Cada identificador deberá ser positivo y único.
    Los campos numéricos deberán rechazar texto y valores fuera del rango
    aplicable; existencias y puntos acumulados no podrán ser negativos.
    \item \textbf{Validación.} Los textos obligatorios no deberán estar vacíos;
    los correos deberán tener formato válido y los teléfonos contener diez
    dígitos. La validación se realizará antes de persistir un cambio.
    \item \textbf{Manejo de errores.} Datos inválidos, identificadores ausentes o
    repetidos y problemas de lectura/escritura no deberán cerrar la aplicación.
    El sistema deberá comunicar mensajes claros y permitir corregir la entrada.
    \item \textbf{Mantenibilidad.} El código deberá conservar la separación por
    capas: interfaz, servicios, repositorios CSV, modelos y validaciones. Cada
    clase y método deberá documentarse con Javadoc.
    \item \textbf{Compatibilidad.} El proyecto deberá compilarse y probarse con
    Maven y Java~25, en UTF-8 y sin dependencias no declaradas en \texttt{pom.xml}.
\end{enumerate}
